\definecolor{appendixblue}{RGB}{43,86,166}
\newcommand{\appendixcontentssection}[3]{%
  \par\noindent
  \hyperref[#3]{{\color{appendixblue}\large\bfseries
   \makebox[2.7em][l]{#1}#2}}%
  \nobreak\hfill
  \hyperref[#3]{{\large\bfseries\pageref*{#3}}}\par
  \vspace{0.55ex}%
}
\newcommand{\appendixcontentssubsection}[3]{%
  \par\noindent\hspace*{2.7em}%
  \hyperref[#3]{{\color{appendixblue}\makebox[3.6em][l]{#1}#2}}%
  \nobreak\leaders\hbox{\kern0.28em.\kern0.28em}\hfill\nobreak
  \hyperref[#3]{\makebox[2.2em][r]{\pageref*{#3}}}\par
  \vspace{0.22ex}%
}

\pdfbookmark[0]{Appendix}{appendix-contents}
\begingroup
\hypersetup{hidelinks}
\begin{center}
  {\color{appendixblue}\fontsize{20}{28}\selectfont\bfseries
   Appendix}
\end{center}
\vspace{2.2em}

\appendixcontentssection{A}{Benchmark Details}{app:benchmark-details}
\appendixcontentssubsection{A.1}{Details of the Evolve Harness}{app:evolve-harness}
\appendixcontentssubsection{A.2}{Details of the Policy Harness}{app:policy-harness}
\appendixcontentssubsection{A.3}{Benchmark Difficulty and Harness Sensitivity}{app:benchmark-item-map}

\vspace{1.4ex}
\appendixcontentssection{B}{Details of Benchmark Construction}{app:stage1-details}
\appendixcontentssubsection{B.1}{Selection of Auxiliary Tasks}{app:auxiliary-task-selection}
\appendixcontentssubsection{B.2}{Selection of Auxiliary Harnesses}{app:auxiliary-harness-selection}

\vspace{1.4ex}
\appendixcontentssection{C}{Experiment Details}{app:experiment-details}
\appendixcontentssubsection{C.1}{Model Configurations and Cost Accounting}{app:model-configurations}
\appendixcontentssubsection{C.2}{Integrity and Reward-Hacking Analysis}{app:integrity-analysis}

\vspace{1.4ex}
\appendixcontentssection{D}{Failure Mode Analysis of Evolution Trajectories}{app:failure-modes}
\appendixcontentssubsection{D.1}{Qwen3.6-27B: Misattributing Regressions to Noise}{app:qwen-failure-mode}
\appendixcontentssubsection{D.2}{DeepSeek-V4-Pro: Premature Plateau and Procedural Evaluations}{app:deepseek-failure-mode}
\appendixcontentssubsection{D.3}{Kimi-K2.7-Code: Careful Rollback but Local Search Saturation}{app:kimi-failure-mode}
\appendixcontentssubsection{D.4}{Cross-Model Failure Patterns}{app:cross-model-failures}

\vspace{1.4ex}
\appendixcontentssection{E}{Prompt Design}{app:prompt-design}
\appendixcontentssubsection{E.1}{Evolver Prompt}{app:evolver-prompt}
\appendixcontentssubsection{E.2}{Seed Policy Prompt}{app:seed-policy-prompt}
\appendixcontentssubsection{E.3}{Judge Prompts}{app:judge-prompts}
\endgroup

\newpage


\section{Benchmark Details}
\label{app:benchmark-details}

This section supplements the benchmark description with implementation and diagnostic details.  It describes the fixed evolve harness, the seed policy harness, and the task-level difficulty and harness-sensitivity characteristics of the constructed benchmark suites.

\subsection{Details of the Evolve Harness}
\label{app:evolve-harness}

The fixed evolve harness $\mathcal{H}_{\mathrm{evo}}$ lets every evolver inspect,
revise, and evaluate a policy harness under the same prompt, tools, state
management, and resource accounting.  Table~\ref{tab:evolve-harness-components}
summarizes its components.

\begin{table}[htbp]
\centering
\small
\caption{Components of the evolve harness.}
\label{tab:evolve-harness-components}
\begin{tabularx}{\textwidth}{p{0.21\textwidth} X}
\toprule
\textbf{Component} & \textbf{Function} \\
\midrule
Evolution orchestrator & Defines the optimization objective, hypothesis-driven
workflow, editable surfaces, integrity rules, and termination protocol. \\
Engineering tools & Provide file inspection and editing, code search, shell
execution, web access, and read-only delegated exploration. \\
Research skills & Support rollout slicing, behavioral comparison, multi-eval
analysis, experiment logging, insight retention, and architecture-level
reflection. \\
Evaluation interface & Evaluates immutable harness snapshots asynchronously and
returns scores, per-domain results, task outcomes, and failure diagnostics. \\
\makecell[l]{Context and state\\management} & Tracks budgets, evaluations, best revisions, and a persistent
experiment ledger, while compacting old interaction history for long runs. \\
\makecell[l]{Sandbox and snapshot\\management} & Separates evolver evidence from policy inputs, records
revisions and trajectories, and freezes one harness for held-out evaluation. \\
\bottomrule
\end{tabularx}
\end{table}

\paratitle{Evolution orchestrator.}
The orchestrator separates the editable policy harness from a scratch
workbench and promotes a diagnose--hypothesize--edit--evaluate cycle over all
harness surfaces.  A live state view reports evaluation, step, time, and token
budgets, recent scores, and experiment history; limits are enforced by the
harness.

\paratitle{Engineering tools.}
File, search, shell, and web tools support implementation and external research.
A read-only subagent can inspect code, validation artifacts, and web sources in
an independent context, but cannot edit or evaluate the harness.

\paratitle{Research skills.}
Six progressively disclosed skills provide rollout filtering
(\texttt{query\_rollouts}), task-paired comparison (\texttt{diff\_task}), and
multi-evaluation analysis (\texttt{analyze\_evals}).  The ledger skills record
hypotheses (\texttt{log\_experiment}), close them from evidence
(\texttt{record\_insight}), and trigger broader redesign after repeated local
failures (\texttt{architecture\_checkpoint}).  All are optional views over raw,
evolver-visible artifacts.

\paratitle{Evaluation interface.}
\texttt{run\_train\_eval} consumes one iteration and asynchronously evaluates an
immutable snapshot; at most one evaluation is active.  \texttt{await\_eval}
returns the headline metric, changes from previous and best revisions,
per-domain results, worst tasks, and failure classes.  Full results, feedback,
logs, and trajectories remain available.

\paratitle{Context and state management.}
The harness records interactions, tool calls, evaluations, resource use, and
snapshots; the ledger persists resolved hypotheses.  Near the context limit,
older messages are summarized while recent turns and all disk artifacts remain
available.

\paratitle{Sandbox and snapshot management.}
The evolver may read validation evidence but write only the policy harness and
workbench.  Policy rollouts receive a read-only snapshot, public task, and
isolated workspace, never answers, scorers, evolver files, or held-out data.
Integrity checks scan source and trajectories; held-out evaluation begins only
after the selected revision is frozen.

\subsection{Details of the Policy Harness}
\label{app:policy-harness}

All runs start from the CodeAct seed policy harness
$H_0$~\cite{wang2024executable} in Figure~\ref{fig:seed-policy-harness}.  The
fixed policy model alternates shell commands and append-only observations until
\texttt{finish}, a tool-used answer, or a resource limit.  Domain tools,
planning, memory, and verification are left for evolution.

\begin{figure}[htbp]
    \centering
    \begin{tikzpicture}[
        node distance=1.35cm and 1.55cm,
        basebox/.style={rounded corners=2pt, minimum height=0.72cm,
                        minimum width=2.45cm, align=center, font=\small},
        neutralbox/.style={basebox, draw=black!55, fill=black!4},
        modelbox/.style={basebox, draw=blue!65, fill=blue!7},
        shellbox/.style={basebox, draw=teal!65, fill=teal!7},
        flow/.style={-{Latex[length=2mm]}, line width=0.8pt, draw=black!65},
        label/.style={font=\scriptsize, text=black!70, fill=white,
                      inner sep=1.5pt}
    ]
        \node[neutralbox] (task) {Public task\\and workspace};
        \node[modelbox, right=of task] (model) {Fixed policy model};
        \node[shellbox, right=of model] (shell) {\texttt{run\_shell\_command}};
        \node[neutralbox, below=0.95cm of model] (answer) {Final answer\\or artifact};

        \draw[flow] (task) -- node[label, above]{instruction} (model);
        \draw[flow] ([yshift=4pt]model.east) --
            node[label, above]{command} ([yshift=4pt]shell.west);
        \draw[flow] ([yshift=-4pt]shell.west) --
            node[label, below]{observation} ([yshift=-4pt]model.east);
        \draw[flow] (model) -- node[label, right]{\texttt{finish}} (answer);
    \end{tikzpicture}
    \caption{The minimal CodeAct policy harness. Shell observations are appended
    to the model context until the task is completed or a resource limit is
    reached.}
    \label{fig:seed-policy-harness}
\end{figure}

\subsection{Benchmark Difficulty and Harness Sensitivity}
\label{app:benchmark-item-map}

Figure~\ref{fig:alpha-beta-item-map} plots harness sensitivity
$\mathrm{Sens}(x)$ and difficulty $1-\mathrm{Perf}(x)$, where
$\mathrm{Perf}(x)$ is mean performance across the twelve auxiliary harnesses.
Selection favors non-negative sensitivity while preserving difficulty coverage,
and the two splits remain aligned.  BrowseComp, HLE, and APEX-Agents skew hard;
GDPval is broader and Claw-Eval easier.  Eight slightly negative Claw-Eval tasks
are quota backfills from sparse strata.

\begin{figure*}[htbp]
    \centering
    \includegraphics[width=\textwidth]{Figures/alpha_beta_item_map.pdf}
    \caption{Difficulty and harness sensitivity of candidate and selected benchmark
    tasks.  Gray, blue, and orange denote unselected, validation, and evaluation
    tasks; marker size denotes coordinate multiplicity.  Annotations give pool
    and selected sizes, and shading marks negative sensitivity.}
    \label{fig:alpha-beta-item-map}
\end{figure*}


\section{Details of Benchmark Construction}
\label{app:stage1-details}
Section~\ref{sec:calibration-tasks} summarizes the two steps of Stage~1: we
first select a 320-task auxiliary suite and then use it to generate and select a
representative set of auxiliary harnesses.  This section provides the full
selection details for both components.  Section~\ref{app:auxiliary-task-selection}
describes the selection of auxiliary tasks, and
Section~\ref{app:auxiliary-harness-selection} describes the selection of
auxiliary harnesses.

\subsection{Selection of Auxiliary Tasks}
\label{app:auxiliary-task-selection}

The first step of Stage~1 selects the auxiliary tasks used to induce harness
diversity.  We
collect Search from MiroRL~\cite{2025mirorl}, RedSearcher~\cite{redsearcher2026},
and internal corpora; Office from anonymized enterprise workflows; and General
from Auto-ClawEval~\cite{li2026clawenvkitautomaticenvironmentgeneration}.  After
validity, rollout-coverage, deduplication, and cross-suite leakage checks, 11,322
candidates remain: 2,988 Search, 7,564 Office, and 770 General.

\paratitle{Rollout-based characterization.}
Using DeepSeek-V4-Flash with the domain-native MiroFlow~\cite{miromind2025mirothinker},
Stirrup~\cite{artificialanalysis2026stirrup}, and Claw-Eval~\cite{claweval}
workflows, we estimate mean score $s(x)$ and turns $\ell(x)$.  We require two
judged Search rollouts and three valid Office/General rollouts.  This process is
independent of the CodeAct seed.

\paratitle{Headroom--horizon selection.}
Within each domain, tasks enter six score bands:
\emph{frontier} ($s=0$), \emph{hard} ($0<s\leq0.10$), \emph{hardish}
($0.10<s\leq0.25$), \emph{mid} ($0.25<s\leq0.50$), \emph{solvable}
($0.50<s\leq0.75$), and \emph{easy} ($s>0.75$), with quota weights
$(0.58,0.24,0.11,0.045,0.02,0.005)$.  Deficits are redistributed, and candidates
within each band are ranked by decreasing $\ell(x)$ with seed 20260716.  This
selects 128 Search, 128 Office, and 64 General tasks, favoring headroom and long
interaction horizons (Figure~\ref{fig:auxiliary-task-selection}).

\begin{figure}[htbp]
    \centering
    \includegraphics[width=1.0\linewidth]{Figures/v3_calibration_selection_paper.pdf}
    \caption{Auxiliary-task selection by headroom and interaction horizon.
    Rows are Search, Office, and General; columns show score bands, turn CDFs,
    and score--turn distributions.  Blue is the candidate pool and red the
    selected 320-task suite.}
    \label{fig:auxiliary-task-selection}
\end{figure}


\subsection{Selection of Auxiliary Harnesses}
\label{app:auxiliary-harness-selection}
After selecting the auxiliary tasks, the second step of Stage~1 uses them to
generate a diverse pool of harness revisions and select a representative
auxiliary harness set.
Claude Opus-4.8, Claude Sonnet-5, GLM-5.2, and GPT-5.6-Sol independently evolve
the same seed CodeAct harness on the 320 auxiliary tasks with
DeepSeek-V4-Flash as the fixed policy model and the main budget (20 iterations,
1,000 steps, 48 hours).  Holding the starting harness, policy model, evolution
protocol, and resource budget constant isolates variation induced by the four
evolvers.  Together, the four evolutions produce 73 evaluated harnesses for the
subsequent selection procedure.

\paratitle{Snapshot deduplication and representation.}
Because successive iterations can produce code-identical harnesses, we first
deduplicate the candidate pool.  Content hashing reduces the 73 evaluated
harnesses to 65 unique revisions; the common seed harness is excluded.  To
compare the remaining revisions on a common basis, we represent each one by its
exercised and coded capabilities, observed tool use, program structure, and
auxiliary performance.  The capability representation covers tooling, context,
control flow, verification, memory, recovery, multi-agent execution, and
routing.

\paratitle{Representative-anchor selection.}
We then select a balanced subset that preserves the progression within each
evolution while covering distinct harness designs.  A deterministic,
final-snapshot-constrained $k$-medoids procedure selects three harnesses per
evolution.  It retains the final revision, preserves capability levels seen in
at least two revisions, minimizes representation error, and breaks ties by
pairwise diversity.  Selecting the same number of harnesses from each evolution
prevents evolutions with more evaluated revisions from dominating the selected
set.  The four evolutions therefore yield 12 harnesses,
$\mathcal{H}_{\mathrm{aux}}$, completing the auxiliary-harness selection step
of Stage~1.  As described in Section~\ref{sec:harness-response-prior},
candidate-task evaluations under this set are subsequently used in Stage~2 to
compute $\mathrm{Sens}(x)$ and $\mathrm{Perf}(x)$. 
% Table~\ref{tab:anchor-pipeline} summarizes the selection outcome.
% \begin{table}[htbp]
% \centering
% \caption{Per-lineage yield of the auxiliary-harness selection pipeline.}
% \label{tab:anchor-pipeline}
% \begin{tabular}{lccc}
% \toprule
% Lineage & Evaluated snapshots & Unique revisions & Selected anchors \\
% \midrule
% Claude Opus-4.8  & 18 & 18 & 3 \\
% Claude Sonnet-5  & 20 & 16 & 3 \\
% GLM-5.2          & 15 & 13 & 3 \\
% GPT-5.6-Sol      & 20 & 18 & 3 \\
% \midrule
% \textbf{Total}   & \textbf{73} & \textbf{65} & \textbf{12} \\
% \bottomrule
% \end{tabular}
% \end{table}


\section{Experiment Details}
\label{app:experiment-details}
This section supplements the experimental setup and reporting protocol.  It documents model configurations and cost accounting, followed by the integrity controls and trajectory audits used to detect and correct reward-hacking behavior.


\subsection{Model Configurations and Cost Accounting}
\label{app:model-configurations}

We evaluate seven frontier and two open-weight evolvers
(Table~\ref{tab:model-configurations}) using thinking mode, the largest supported
context, and the highest exposed reasoning effort.  Requests allow 65,536
output tokens and time out after 600 seconds; default
temperature is $1.0$ unless endpoint-controlled.  Policy and judge are fixed.


\begin{table}[htbp]
\centering
\small
\caption{Model configurations used in the main experiments.  ``Enabled''
denotes a model-specific thinking switch rather than a scalar reasoning-effort
setting.  ``Default'' indicates endpoint-controlled temperature.}
\label{tab:model-configurations}
\begin{tabular*}{\textwidth}{@{\extracolsep{\fill}}lllcc@{}}
\toprule
\textbf{Role} & \textbf{Model} & \textbf{Reasoning} & \textbf{Temp.} & \textbf{Context} \\
\midrule
Evolver & Qwen3.7-Max       & \texttt{xhigh} & 1.0 & 1M \\
Evolver & MiniMax M3        & \texttt{max}   & 1.0 & 1M \\
Evolver & DeepSeek-V4-Pro   & \texttt{max}   & 1.0 & 1M \\
Evolver & Kimi-K2.7-Code    & Enabled        & Default & 256K \\
Evolver & GLM-5.2           & \texttt{max}   & 1.0 & 1M \\
Evolver & GPT-5.6-Sol       & \texttt{max}   & 1.0 & 384K \\
Evolver & Claude Opus-4.8   & \texttt{max}   & 1.0 & 1M \\
Evolver & Qwen3.6-27B       & \texttt{xhigh} & 1.0 & 256K \\
Evolver & Gemma 4 31B       & Enabled        & 1.0 & 256K \\
\midrule
Policy (fixed) & DeepSeek-V4-Flash & \texttt{max} & 1.0 & 256K \\
Judge (fixed)  & Qwen3.7-Plus      & Standard & 0.0 & 1M \\
\bottomrule
\end{tabular*}
\end{table}
\paratitle{Fixed policy and judge.}
DeepSeek-V4-Flash uses maximum reasoning, 256K context, the same output cap, and
a 600-second timeout.  Qwen3.7-Plus grades model-judged tasks at temperature
$0.0$, 1M context, the same cap, and a 1,200-second timeout.  Claude Opus-4.8
uses Anthropic messages; other models use OpenAI-compatible interfaces.

\paratitle{Cost accounting.}
From per-call usage, model $m$ costs
\begin{equation*}
C_m=\frac{1}{10^6}\sum_r
\left(n^{\mathrm{in}}_r p^{\mathrm{in}}_m
+n^{\mathrm{cr}}_r p^{\mathrm{cr}}_m
+n^{\mathrm{cw}}_r p^{\mathrm{cw}}_m
+n^{\mathrm{out}}_r p^{\mathrm{out}}_m\right),
\end{equation*}
where $n$ counts uncached input, cache-read, cache-write, and output tokens, and
$p$ is the corresponding public non-batch list price per million tokens on July
10, 2026.  Request length selects tiered prices.  The cost--performance analysis
includes main-loop, compaction, and subagent calls but excludes policy and judge
costs; full decompositions are retained.  Missing cache counters are charged as
uncached input.

\subsection{Integrity and Reward-Hacking Analysis}
\label{app:integrity-analysis}

Harness evolution deliberately gives the evolver broad autonomy within an independent sandbox.  It can inspect rollout trajectories, task provenance, expected answers, and validation scores; use shell and web tools; and modify the policy harness.  This access is necessary for diagnosing failures and developing general harness improvements, but it also creates substantial opportunities for reward hacking.  During our experiments, we focused on three main forms of hacking behavior and introduced a corresponding control for each one.

\paratitle{(1) Leakage across the evolver--policy boundary.}
An evolver could misuse validation access by transferring protected information into policy execution or by attempting to interfere with the evaluation protocol.  We control this risk through strict sandbox separation.  The evolver may read validation artifacts but can write only to its policy harness and workbench.  Each policy rollout instead receives an immutable harness snapshot, a public task view with answers and scorer fields removed, and a fresh workspace; validation artifacts, evolver notes, sibling workspaces, and held-out data are not mounted.  Scorers, aggregation logic, suite definitions, and iteration accounting remain in the fixed host-side framework.  Every formal validation evaluation consumes one iteration, and held-out evaluation starts only after the selected harness is frozen and returns no evidence to the evolver.

\paratitle{(2) Retrieval of benchmark answers.}
A model with web access may search for benchmark-specific answers, answer keys, or evaluation artifacts instead of solving a task from permitted evidence. Simply disabling web access is not appropriate because retrieval is intrinsic to Search tasks.  We therefore scan content returned by web and shell tools for known patterns associated with benchmark datasets, answer keys, reference artifacts, and evaluation dumps.  The check applies to information delivered to the policy rather than to attempted queries: an unsuccessful request is not penalized, and a request blocked before protected content reaches the policy is not treated as a violation.  For a confirmed match, we record the source and tool event, mark the trial as a policy violation, and set all native score fields to zero before computing domain scores and Overall Score.

\paratitle{(3) Evasion of the retrieval safeguard.}
A model may also attempt to bypass the preceding scanner by paraphrasing, encoding, obfuscating, or concealing retrieved information.  Because pattern matching alone cannot reliably detect such behavior, we apply the same Codex-based semantic audit to the complete evolver and policy trajectories of every evolution run.  The audit checks whether protected information reached policy execution, whether retrieval targeted benchmark-specific artifacts, and whether behavior attempted to conceal detector-visible evidence.  Among all evaluated models, only MiniMax M3 exhibited this detector-evasion behavior.  The audit identified the affected trials; we set their scores to zero, recomputed the aggregates, and report the corrected MiniMax M3 result.  Tool calls, observations, harness snapshots, and violation records are retained for inspection.


\section{Failure Mode Analysis of Evolution Trajectories}
\label{app:failure-modes}

We present the complete evolution process of three lower-scoring models and analyze the failure modes that limited their final performance.  The analysis draws on transcripts, experiment ledgers, validation digests, and harness snapshots.  Scores are validation Overall Scores on the same 0--100 scale used in the main experiments; I$i$ denotes iteration $i$, and evaluation-suite results are excluded.

\begin{figure}[htbp]
    \centering
    \includegraphics[width=\linewidth]{Figures/failure_mode_trajectories.pdf}
    \caption{Validation Overall Score trajectories.  Bands show intervention
    surfaces; stars and squares mark best and frozen revisions.  Catastrophic
    revisions are annotated.  Kimi restored I13 after I20.}
    \label{fig:failure-mode-trajectories}
\end{figure}


\begin{table}[htbp]
\centering
\caption{Research-budget use and revision selection.  ``Rollout accesses''
counts raw-artifact accesses; Kimi restored the byte-identical I13 snapshot
after I20.}
\small
\begin{tabular*}{\linewidth}{@{\extracolsep{\fill}}lcccc@{}}
\toprule
Evolver & Iterations & Steps & Rollout accesses & Best $\rightarrow$ frozen \\
\midrule
Qwen3.6-27B      & 18/20 & 349/1000 & 16 & 49.7 (I10) $\rightarrow$ 45.4 (I18) \\
DeepSeek-V4-Pro  & 15/20 & 202/1000 & 4  & 46.5 (I3) $\rightarrow$ 42.6 (I15) \\
Kimi-K2.7-Code   & 20/20 & 549/1000 & 52 & 45.9 (I13) $\rightarrow$ 45.9 (I13) \\
\bottomrule
\end{tabular*}
\label{tab:failure-mode-budget}
\end{table}
\subsection{Qwen3.6-27B: Misattributing Regressions to Noise}
\label{app:qwen-failure-mode}

Qwen rapidly added retrieval, file/Python tools, routing, integrity filters,
verification, and recovery, rising from 33.0 to 49.7 at I10.  Later it bundled
context and verifier changes; a malformed I15 verifier crashed 159/160 tasks
and reduced the score to 0.3, a defect missed by local preflight.  Repairs
recovered only to 45.4 at I18.

Despite a 2.2-point Overall Score range across byte-identical I8/I10/I12
revisions, Qwen attributed the 4.3-point I10--I18 gap largely to noise without
replicating or restoring I10/I12.
It froze I18 with two iterations and 651 steps unused.  The failure combines
bundled edits, missing smoke checks, and weak best-snapshot management.

\subsection{DeepSeek-V4-Pro: Premature Plateau and Procedural Evaluations}
\label{app:deepseek-failure-mode}

DeepSeek reached 46.5 at I3 after adding search, fetch, file/Python tools, URL
filters, and domain guidance.  It then tuned prompts, filters, and completion
rules while accessing raw rollouts only four times; I4--I10 ranged from 35.9 to
43.0.  Requiring three searches at I11 reduced the score to 42.5 and tripled
runtime; I12 reflection reached 40.6.

After an attempted early stop, I13--I15 scored 40.8, 40.4, and 42.6, with the
last two evaluating unchanged code.  DeepSeek labeled I3 an outlier without
reevaluating it and froze I15 with five iterations and 798 steps unused.  This
is aggregate hill climbing without causal localization, followed by procedural
budget use and premature plateau declaration.

\subsection{Kimi-K2.7-Code: Careful Rollback but Local Search Saturation}
\label{app:kimi-failure-mode}

Kimi recovered from an invalid I1 architecture (score 2.1) to 45.1 at I6,
tested and reverted regressive context and finish controls, then reached 45.9
at I13 with a step-budget planner.  Seven further tests of planning, APEX
instructions, thinking, checklists, step limits, and verification did not
improve it.  Kimi logged these results, restored byte-identical I13, and froze
the best revision.

Its limitation was search efficiency: despite exhausting all 20 evaluation
iterations and accessing 52 rollouts, it used only 549/1,000 steps.  Post-I13
work remained local prompt, gate, context, and threshold variants rather than
deeper planning or verification redesign.  Rollback discipline prevented
regression but not architectural saturation.

\subsection{Cross-Model Failure Patterns}
\label{app:cross-model-failures}

All three add obvious capabilities quickly but later favor local prompts,
filters, thresholds, and gates over causal, task-paired redesign.  Qwen and
DeepSeek stop early; Kimi exhausts evaluations on low-yield search.  Only Kimi
restores its best snapshot.  Stronger evolution therefore needs paired failure
analysis, one falsifiable mechanism per iteration, cheap preflight checks,
automatic best-revision recovery, and an architectural-reset trigger after
repeated local failures.


\section{Prompt Design}
\label{app:prompt-design}

This section summarizes the prompts used throughout the Evo-Bench evaluation process.  They comprise a fixed evolver prompt, an editable policy prompt, and fixed evaluation-side judge prompts.  Evolvers may modify only the policy layer; policies never receive references, rubrics, or judge instructions.


\subsection{Evolver Prompt}
\label{app:evolver-prompt}

All nine evolvers receive the same template below, apart from serialization.
Bracketed fields supply domains, metric, baseline, and budget; only repeated
motivational text and expanded schemas/catalogs are omitted.

\begin{tcolorbox}[
    breakable,
    colback=black!2,
    colframe=black!35,
    boxrule=0.5pt,
    arc=1mm,
    left=1.5mm,
    right=1.5mm,
    top=1mm,
    bottom=1mm,
    title={Fixed evolver prompt template},
    fonttitle=\bfseries\small]
\footnotesize
\textbf{Identity and objective.}
You are a harness research engineer operating in an isolated sandbox.  Improve
the editable policy harness to maximize the held-out headline metric.  Use
validation only, prefer general mechanisms, and diagnose trajectories before
structural changes; treat regressions as evidence.

\textbf{Boundaries and integrity.}
The policy harness at \texttt{[POLICY]} is the product that will be frozen and
scored; \texttt{[WORKBENCH]} is unscored scratch space.  Validation answers,
feedback, and rollouts are diagnostic evidence, not policy resources.  Do not
modify evaluation, transfer protected evidence, exploit artifacts, or conceal
policy actions.

\textbf{Tools.}
Use engineering tools and read-only subagents for scoped research;
\texttt{get\_research\_state} for budgets;
\texttt{run\_train\_eval}/\texttt{await\_eval} for formal evaluation; and
research skills for rollout analysis and the experiment ledger.  Call
\texttt{complete\_evolution} only after restoring the selected revision.

\textbf{Experimental workflow.}
Treat the clean seed as iteration 0.  Repeatedly inspect failures, state a
falsifiable mechanism, edit and locally test, evaluate an immutable snapshot,
and retain or revert from evidence.  Any harness surface may be redesigned.
Analyze while evaluations run and restore the best revision before freezing.

\textbf{Opening run specification.}
The run contains \texttt{[DOMAIN DESCRIPTIONS]} and optimizes
\texttt{[HEADLINE METRIC]}.  The clean baseline is
\texttt{[BASELINE DIGEST]}.  Redesign the minimal seed within the evaluation,
step, and time budgets; evaluation-suite tasks are unavailable.

\textbf{Research state.}
\texttt{iteration [USED]/[MAX]; step [USED]/[MAX]; [TIME] remaining;
tokens [COUNT]; current/best score [VALUES]; active evaluation [STATE].}
\end{tcolorbox}

The dataset block gives task families and $2{:}2{:}1$ aggregation, but no
held-out instances.  Live state comes from tools rather than self-report.  Near
the context limit, a compaction prompt preserves per-iteration mechanisms,
scores, hypotheses, best snapshot, open questions, and paths; only old messages
are replaced, not disk artifacts.

\subsection{Seed Policy Prompt}
\label{app:seed-policy-prompt}

The minimal seed specifies only CodeAct interaction and workspace boundaries.

\begin{tcolorbox}[
    breakable,
    colback=blue!2,
    colframe=blue!35,
    boxrule=0.5pt,
    arc=1mm,
    left=1.5mm,
    right=1.5mm,
    top=1mm,
    bottom=1mm,
    title={Seed policy system prompt},
    fonttitle=\bfseries\small]
{\footnotesize\ttfamily\raggedright
You are a CodeAct-style task-solving agent.\par
\medskip
You work inside a dedicated task workspace directory. Every
run\_shell\_command runs with that workspace as its working directory, so
relative paths resolve against it. Create, edit, and read your files there.
The concrete absolute workspace path is given in the task message. Use
run\_shell\_command to inspect files, run programs, edit files, and verify your
work. Use finish when the task is done.\par
\medskip
Rules:\par
- Keep all files inside the task workspace.\par
- Do not access paths outside it unless explicitly shown in the task.\par
- Prefer simple shell commands and small file edits.\par
- If a command fails, inspect the output and continue.\par
- Call finish with a concise final answer or artifact path.\par}
\end{tcolorbox}

For task $x$, the seed task analyzer constructs the user message as follows:

\begin{tcolorbox}[
    colback=black!2,
    colframe=black!30,
    boxrule=0.5pt,
    arc=1mm,
    left=1.5mm,
    right=1.5mm,
    top=1mm,
    bottom=1mm]
{\footnotesize\ttfamily\raggedright
Task id: [TASK\_ID]\par
Task workspace (your working directory): [WORKSPACE]\par
\medskip
[PUBLIC TASK PROMPT]\par
\medskip
Public files:\par
- [WORKSPACE-RELATIVE FILE]\par}
\end{tcolorbox}

The file block is optional and the seed omits the domain label.  Scorer fields,
expected outputs, anchors, and rationales are stripped.  Evolved prompt and code
are frozen together.

\subsection{Judge Prompts}
\label{app:judge-prompts}

Judge instructions are fixed and post-rollout.  Deterministic scorers use no
prompt; model judges use Section~\ref{app:model-configurations} and receive
references only inside evaluation (Table~\ref{tab:judge-prompt-design}).

\begin{table}[htbp]
\centering
\small
\caption{Fixed judge-prompt interfaces.  All structured outputs are parsed and
retried on formatting failure; failed grading runs are surfaced as evaluation
errors rather than silently treated as valid model failures.}
\label{tab:judge-prompt-design}
\begin{tabularx}{\textwidth}{p{0.18\textwidth} p{0.34\textwidth} X}
\toprule
\textbf{Scorer} & \textbf{Judge evidence} & \textbf{Required decision} \\
\midrule
BrowseComp / semantic equivalence & Question, protected short reference answer,
and policy final answer & \texttt{CORRECT: YES/NO} plus one-sentence reason;
formatting and equivalent phrasing are ignored. \\

HLE & Question, policy response, and precise protected answer & Extract the
final answer, judge strict equivalence (with a small numerical tolerance), and
return reasoning, a yes/no verdict, and confidence. \\

GDPval absolute & Task, weighted positive/penalty criteria, parsed deliverable
text, formulas, and rendered pages/slides & JSON with points and evidence-backed
reason for every criterion, plus an overall summary. \\

GDPval pairwise & The same rubric plus policy and frozen-seed deliverables,
including text and images & For every criterion choose A, B, or tie.  Both A/B
orders are judged and averaged to reduce position bias. \\

Claw-Eval & Complete executed tool-use trace, task-private grader, rubric, and
reference state & Benchmark-native completion, robustness, communication, and
safety scores, combined by the official grader. \\

APEX-Agents & Policy trajectory, environment-state difference, artifacts, and
task-private verifiers inside the APEX sandbox & Per-verifier rationale and
score; a task passes only when every required criterion passes. \\
\bottomrule
\end{tabularx}
\end{table}

Semantic grading ignores surface form; HLE uses extract-then-compare.  GDPval
uses rendered images for visual criteria and parsed text for facts/formulas;
pairwise grading averages both A/B orders.  APEX-Agents and Claw-Eval retain
native graders.  Validation rationales are evolver-visible but policy-hidden;
held-out feedback is generated only after freezing.
